\section{Status and the remaining route to a proof}\label{sec:outlook}

The combined investigation remains open as a field-theoretic research
problem. The original fixed-scalar family has an exact Loewner
description but little contour freedom at fixed charge. Tangent
coupling removes that geometric restriction and admits the full
chiral defect projectors. Its ordinary same-charge scalar pairing
vanishes. Physical reflection produces a different, nonzero free
pairing and a smooth reference join, but it changes the branch
charges and has no established protection theorem. These conclusions
can coexist: they concern different pairings of the same basic fields.

\subsection{Statement ledger}

\renewcommand{\arraystretch}{1.15}
\begin{longtable}{@{}p{0.36\textwidth}p{0.59\textwidth}@{}}
\toprule
Statement & Status in this draft\\
\midrule
\endfirsthead
\toprule
Statement & Status in this draft (continued)\\
\midrule
\endhead
\bottomrule
\endfoot
Weil positivity criterion & Classical background theorem; target
normalization given in \eqref{eq:weil}.\\
Beta/comb and rational factorization & Exact identities reproduced
with their convergence region; no contraction conclusion.\\
Free even angular kernel & Positive for $\rho<1$; its difference
energy has a positive limit. The opposite contact subtraction does not.\\
Auxiliary gamma partition ratio & Exact finite Gaussian product and
renormalized limit; no physical mode selection or transfer identified.\\
Constant-driver dictionary and rigidity & Exact contour identity and
proof under the specified fixed-scalar bulk charge.\\
Smooth transport and inverse-map variation & Exact matrix-field ODE
identities away from singularities; no rough-trace quantum construction.\\
Tangent-map defect supersymmetry & Complex chiral projector ranks proved
for the real constant-map ansatz.\\
Same-charge scalar endpoint pairing & Zero free contraction; vanishing
expectation conditional on an invariant vacuum and operator definition.\\
Ordinary reflected network & Exact classical map and color gluing;
free endpoint positivity proved. Interacting positivity is conditional.\\
Rotation and odd-H tests & Explicit failures on the stated free
algebras, with ordinary-adjoint positive controls.\\
Smooth join and bulge response & Direct Feynman-gauge Gaussian bulk
sector computed; full interacting response remains open.\\
Arithmetic identification & No prime-weight mechanism, complete
transfer, or proof of positivity of $Q_{0,L}$ has been obtained.\\
\end{longtable}
\renewcommand{\arraystretch}{1}

\subsection{The next discriminating calculation}

The immediate continuation is the complete first interaction-order
shape response of \eqref{eq:network}, beginning with a positive bulge
$\eta_0>0$ in \eqref{eq:bulge}. Fix the gauge, endpoint normalization,
reference color sector and renormalization prescription before
comparing deformations. Alongside the direct bulk exchange, include
endpoint-to-line gauge exchange, the defect H-scalar derivative
coupling, endpoint self-energy, and all defect and junction
counterterms at that order. The interaction vertices must be taken
from a consistent defect action, such as the conventions used in
\cite{DFO,Baker}. Endpoint renormalization may cancel in a specified
ratio, but that cancellation must be shown.

The flat limit needs separate analysis. At $\eta=0$ the entire
contour lies on the defect; for $\eta>0$ only its endpoints do.
Defect contact terms can therefore make the full small-height
expansion nonuniform. Analyticity of the bulk integral in $\eta^2$
does not justify analyticity of the complete answer: terms linear
in height or involving logarithms have not been excluded. A complete
calculation at positive height would first determine whether the
nonzero bulk response survives, cancels, or requires a different
operator definition.

\subsection{What would turn the construction into an RH route}

A precise sufficient endpoint would be a Hilbert space $\mathcal H_L$
and a linear map $J_L$ defined for every compactly supported smooth
$f$ such that
\begin{equation}\label{eq:proof-target}
 Q_{0,L}[f]=\norm{J_Lf}_{\mathcal H_L}^2
 \qquad\text{for every }L>0.
\end{equation}
Alternatively, an independently defined positive system could realize
the causal transfer \eqref{eq:transfer} with
$\norm{V_{\omega,L}}\leq1$ for every $L$ and shifts tending to zero.
Equation~\eqref{eq:first-variation-arithmetic} would then give the
same positivity conclusion. Either identity must be derived, rather
than built into an assumed positive metric.

For the present proposal this entails the following concrete steps.
\begin{enumerate}
 \item Define a nonzero observable and its physical adjoint, including
 the reference sector, all endpoint and junction data, a regulator,
 and a limit preserving the required positivity.
 \item Derive an evolution law for the relevant averaged observable.
 The insertions in Section~\ref{sec:variation} must close in a
 specified sector or be controlled by a proved identity. A common
 bulk charge alone does not establish this.
 \item Identify the contour or boundary parameter with the arithmetic
 shift and frequency. The capacity relation
 \eqref{eq:capacity-dictionary} and the auxiliary Gaussian ratio
 \eqref{eq:gaussian-gamma} are partial dictionaries with different
 meanings; neither makes this identification.
 \item Produce the full arithmetic form: the local normalization
 $w_0$, the pole contribution, and each
 $\Lambda(n)/\sqrt n$ translation with its sign. Equivalently produce
 the full transfer, including the prime comb and rational correction.
 A free spectral resemblance or a positive kernel cannot replace
 these exact terms.
 \item Prove the equality and bounds for all test functions and an
 unbounded range of interval lengths, with the limits needed in
 \eqref{eq:proof-target} or \eqref{eq:first-variation-arithmetic}.
 Finite Gram matrices and perturbative coefficients are diagnostics
 for this task, not substitutes for it.
\end{enumerate}

The present work supplies explicit objects on which to test the first
two steps and exposes where protection and positivity separate. Its
value for the RH program depends on making one of those tests succeed
and then deriving the missing arithmetic structure.
